%
% Template Laporan Skripsi/Thesis 
%
% @author  Andreas Febrian, Lia Sadita 
% @version 1.03
%
% Dokumen ini dibuat berdasarkan standar IEEE dalam membuat class untuk 
% LaTeX dan konfigurasi LaTeX yang digunakan Fahrurrozi Rahman ketika 
% membuat laporan skripsi. Konfigurasi yang lama telah disesuaikan dengan 
% aturan penulisan thesis yang dikeluarkan UI pada tahun 2008.
%

%
% Tipe dokumen adalah report dengan satu kolom. 
%
\documentclass[12pt, a4paper, onecolumn, oneside, final]{report}

% Load konfigurasi LaTeX untuk tipe laporan thesis
\usepackage{_internals/uithesis}

% Daftar pemenggalan suku kata dan istilah dalam LaTeX
%
% Hyphenation untuk Indonesia 
%
% @author  Andreas Febrian
% @version 1.00
% 
% Tambahkan cara pemenggalan kata-kata yang salah dipenggal secara otomatis 
% oleh LaTeX. Jika kata tersebut dapat dipenggal dengan benar, maka tidak 
% perlu ditambahkan dalam berkas ini. Tanda pemenggalan kata menggunakan 
% tanda '-'; contoh:
% menarik
%   --> pemenggalan: me-na-rik
%

\hyphenation{
    % alphabhet A
    a-na-li-sa a-tur 
    a-pli-ka-si 
    % alphabhet B
    ba-ngun-an 
    be-be-ra-pa 
    ber-ge-rak
    ber-ke-lan-jut-an 
    ber-pe-nga-ruh 
    % alphabhet C
    ca-ri
    % alphabhet D
    di-sim-pan di-pim-pin de-ngan da-e-rah di-ba-ngun da-pat di-nya-ta-kan 
    di-sim-bol-kan di-pi-lih di-li-hat de-fi-ni-si
    % alphabhet E
    e-ner-gi eks-klu-sif
    % alphabhet F
    fa-si-li-tas
    % alphabhet G
    ga-bung-an ge-rak
    % alphabhet H
    ha-lang-an
    % alphabhet I
    % alphabhet J
    % alphabhet K
    ke-hi-lang-an
    ku-ning 
    kua-li-tas ka-me-ra ke-mung-kin-an ke-se-pa-ham-an
    % alphabhet L
    ling-kung-an
    % alphabhet M
    me-neng-ah
    meng-a-tas-i me-mung-kin-kan me-nge-na-i me-ngi-rim-kan 
    meng-u-bah meng-a-dap-ta-si me-nya-ta-kan mo-di-fi-ka-si
    meng-a-tur
    % alphabhet N
    nya-ta non-eks-klu-sif
    % alphabhet O
    % alphabhet P
	pe-nye-rap-an 
	pe-ngon-trol
    pe-mo-del-an
    pe-ran  pe-ran-an-nya
    pem-ba-ngun-an pre-si-den pe-me-rin-tah prio-ri-tas peng-am-bil-an 
    peng-ga-bung-an pe-nga-was-an pe-ngem-bang-an 
    pe-nga-ruh pa-ra-lel-is-me per-hi-tung-an per-ma-sa-lah-an 
    pen-ca-ri-an peng-struk-tur-an
    % alphabhet Q
    % alphabhet R
    ran-cang-an
    % alphabhet S
    si-mu-la-si sa-ngat
    % alphabhet T
    te-ngah
    ter-da-pat
    % alphabhet U
    % alphabhet V
    % alphabhet W
    % alphabhet X
    % alphabhet Y
    % alphabhet Z
    % special
}

% Load konfigurasi khusus untuk laporan yang sedang dibuat
%-----------------------------------------------------------------------------%
% Informasi Mengenai Dokumen
%-----------------------------------------------------------------------------%
% 
% Judul laporan. 
\var{\judul}{Judul Skripsi/Thesis/Disertasi}
% 
% Tulis kembali judul laporan, kali ini akan diubah menjadi huruf kapital
\Var{\Judul}{Judul Skripsi/Thesis/Disertasi}
% 
% Tulis kembali judul laporan namun dengan bahasa Ingris
\var{\judulInggris}{Unknown Title for Final Report/Thesis/Disertation}

% 
% Tipe laporan, dapat berisi Skripsi, Tugas Akhir, Thesis, atau Disertasi
\var{\type}{Skripsi}
% 
% Tulis kembali tipe laporan, kali ini akan diubah menjadi huruf kapital
\Var{\Type}{Skripsi}
% 
% Tulis nama penulis 
\var{\penulis}{Nama Penulis}
% 
% Tulis kembali nama penulis, kali ini akan diubah menjadi huruf kapital
\Var{\Penulis}{Nama Penulis}
% 
% Tulis NPM penulis
\var{\npm}{NPM}
% 
% Tuliskan Fakultas dimana penulis berada
\Var{\Fakultas}{Ilmu Komputer}
\var{\fakultas}{Ilmu Komputer}
% 
% Tuliskan Program Studi yang diambil penulis
\Var{\Program}{ILMU KOMPUTER}
\var{\program}{Ilmu Komputer}
% 
% Tuliskan tahun publikasi laporan
\Var{\bulan}{Juli}
\Var{\tahun}{2019}
% 
% Tuliskan gelar yang akan diperoleh dengan menyerahkan laporan ini
\var{\gelar}{Sarjana Ilmu Komputer}
% 
% Tuliskan tanggal pengesahan laporan, waktu dimana laporan diserahkan ke 
% penguji/sekretariat
\var{\tanggalPengesahan}{XX Juli 2019} 
% 
% Tuliskan tanggal keputusan sidang dikeluarkan dan penulis dinyatakan 
% lulus/tidak lulus
\var{\tanggalLulus}{XX Juli 2019}
% 
% Tuliskan pembimbing 
\var{\pembimbing}{Prof. ???}
% 
% Alias untuk memudahkan alur penulisan paa saat menulis laporan
\var{\saya}{Penulis}

%-----------------------------------------------------------------------------%
% Judul Setiap Bab
%-----------------------------------------------------------------------------%
% 
% Berikut ada judul-judul setiap bab. 
% Silahkan diubah sesuai dengan kebutuhan. 
% 
\Var{\kataPengantar}{Kata Pengantar}
\Var{\babSatu}{Pendahuluan}
\Var{\babDua}{Sekilas Mengenai \latex}
\Var{\babTiga}{Notasi Matematik}
\Var{\babEmpat}{Struktur Berkas}
\Var{\babLima}{Perintah dalam uithesis.sty}
\Var{\babEnam}{Bab Enam}
\Var{\kesimpulan}{Kesimpulan dan Saran}

% Daftar istilah yang mungkin perlu ditandai 
%
% @author  Andreas Febrian
% @version 1.00
% 
% Mendaftar seluruh istilah yang mungkin akan perlu dijadikan 
% italic atau bold pada setiap kemunculannya dalam dokumen. 
% 

\var{\license}{\f{Creative Common License 1.0 Generic}}
\var{\bslash}{$\setminus$}

% Awal bagian penulisan laporan
\begin{document}
%
% Sampul Laporan
%
% Sampul Laporan

%
% @author  unknown
% @version 1.01
% @edit by Andreas Febrian
%

\begin{titlepage}
    \begin{center}    
        \begin{figure}
            \begin{center}
                \includegraphics[width=2.5cm]{_internals/makara.eps}
            \end{center}
        \end{figure}    
        \vspace*{0cm}
        \bo{
        	UNIVERSITAS INDONESIA\\
        }
        
        \vspace*{1.0cm}
        % judul thesis harus dalam 14pt Times New Roman
        \bo{\Judul} \\[1.0cm]

        \vspace*{2.5 cm}    
        % harus dalam 14pt Times New Roman
        \bo{\Type}

        \vspace*{3 cm}       
        % penulis dan npm
        \bo{\Penulis} \\
        \bo{\npm} \\

        \vspace*{5.0cm}

        % informasi mengenai fakultas dan program studi
        \bo{
        	FAKULTAS \Fakultas\\
        	PROGRAM STUDI \Program \\
        	DEPOK \\
        	\tahun
        }
    \end{center}
\end{titlepage}


%
% Gunakan penomeran romawi
\pagenumbering{roman}

%
% load halaman judul dalam
\addChapter{HALAMAN JUDUL}
%
% Halaman Judul Laporan 
%
% @author  unknown
% @version 1.01
% @edit by Andreas Febrian
%

\begin{titlepage}
    \begin{center}\begin{figure}
            \begin{center}
                \includegraphics[width=2.5cm]{pics/makara.png}
            \end{center}
        \end{figure}    
        \vspace*{0cm}
        \bo{
        	UNIVERSITAS INDONESIA\\
        }
        
        \vspace*{1.0cm}
        % judul thesis harus dalam 14pt Times New Roman
        \bo{\Judul} \\[1.0cm]

        \vspace*{2.5 cm}    
        % harus dalam 14pt Times New Roman
        \bo{\Type} \\
        % keterangan prasyarat
        \bo{Diajukan sebagai salah satu syarat untuk memperoleh gelar \\
        \gelar}\\

        \vspace*{3 cm}       
        % penulis dan npm
        \bo{\Penulis} \\
        \bo{\npm} \\

        \vspace*{5.0cm}

        % informasi mengenai fakultas dan program studi
        \bo{
        	FAKULTAS \Fakultas\\
        	PROGRAM STUDI \Program \\
        	DEPOK \\
        	\bulanTahun
        }
    \end{center}
\end{titlepage}

%
% setelah bagian ini, halaman dihitung sebagai halaman ke 2
\setcounter{page}{2}

%
% load halaman pengesahan
\addChapter{LEMBAR PERSETUJUAN}
%
% Halaman Pengesahan
%
% @author  Andreas Febrian
% @version 1.01
%

\chapter*{HALAMAN PERSETUJUAN}

\vspace*{0.2cm}
\noindent 

\noindent
\begin{tabular}{l l p{11cm}}
	\bo{Judul}&: & \judul \\ 
	\bo{Nama}&: & \penulis \\
	\bo{NPM}&: & \npm \\
\end{tabular} \\

\vspace*{1.2cm}

\noindent Laporan \type~ini telah diperiksa dan disetujui.\\[0.3cm]
\begin{center}
\tanggalPengesahan \\[2cm]


\underline{\pembimbing}\\[0.1cm]
Pembimbing \type
\end{center}

\newpage
%
% load halaman orisinalitas 
\addChapter{LEMBAR PERNYATAAN ORISINALITAS}
\input{src/00-front_matter/02-orisinalitas}
%
%
\addChapter{LEMBAR PENGESAHAN}
\input{src/00-front_matter/03-pengesahan_sidang}
%
%
\addChapter{\kataPengantar}
%-----------------------------------------------------------------------------%
\chapter*{Kata Pengantar}
%-----------------------------------------------------------------------------%
Template ini disediakan untuk orang-orang yang berencana menggunakan 
\latex~untuk membuat dokumen tugas akhirnya. 
Mengapa \latex? 
Ada banyak hal mengapa menggunakan \latex, diantaranya:

\begin{enumerate}
	\item \latex~membuat kita jadi lebih fokus terhadap isi dokumen, bukan 
		tampilan atau halaman. 
	\item \latex~memudahkan dalam penulisan persamaan matematis. 
	\item Adanya automatis dalam penomoran caption, bab, subbab, subsubbab, 
		referensi, dan rumus. 
	\item Adanya automatisasi dalam pembuatan daftar isi, daftar gambar, dan
		daftar tabel. 
	\item Adanya kemudahan dalam memberikan referensi dalam tulisan dengan 
		menggunakan label. Cara ini dapat meminimalkan kesalahan pemberian 
		referensi. 
\end{enumerate}

Template ini bebas digunakan dan 
didistribusikan sesuai dengan aturan \license, yang secara sederhana berisi: 

\begin{figure}
	\centering
	\includegraphics[width=0.74\textwidth]
		{assets/pics/creative_common.png}
	\caption{\license}
	\label{fig:lisensi}
\end{figure}

\pic~\ref{fig:lisensi} diambil dari 
\url{http://creativecommons.org/licenses/by-nc-sa/1.0/deed.en_CA}. 
Jika ingin mengentahui lebih lengkap mengenai \license, silahkan buka 
\url{http://creativecommons.org/licenses/by-nc-sa/1.0/legalcode}. 
Seluruh dokumen yang dibuat dengan menggunakan template ini sepenuhnya 
menjadi hak milik pembuat dokumen dan bebas didistribusikan sesuai dengan 
keperluan masing-masing. 
Lisensi hanya berlaku jika ada orang yang membuat template baru dengan 
menggunakan template ini sebagai dasarnya. 

Dokumen ini dibuat dengan \latex~juga. Untuk meyakinkan Anda, coba lihat 
properti dari dokumen ini dan Anda akan menemukan bagian seperti 
\pic~\ref{fig:pdflatex}. 
Dokumen ini dimaksudkan untuk memberikan gambaran kepada Anda seperti apa 
mudahnya menggunakan \latex~dan juga memperlihatkan betapa bagus dokumen 
yang dihasilkan. 
Seluruh url yang Anda temukan dapat Anda klik. 
Seluruh referensi yang ada juga dapat diklik. 
Untuk mengerti template yang disediakan, Anda tetap harus membuka kode 
\latex~dan bermain-main dengannya. 
Penjelasan dalam PDF ini masih bersifat gambaran dan tidak begitu 
mendetail, dapat dianggap sebagai pengantar singkat. 
Jika Anda merasa kesulitan dengan template ini, mungkin ada baiknya 
Anda belajar sedikit dasar-dasar \latex. 

\begin{figure}
	\centering
	\includegraphics[width=0.54\textwidth]
		{assets/pics/mark.png}
	\caption{Dokumen Dibuat dengan PDFLatex}
	\label{fig:pdflatex}
\end{figure}

Semoga template ini dapat membantu orang-orang yang ingin mencoba menggunakan 
\latex. Semoga template ini juga tidak berhenti disini dengan ada kontribusi 
dari para penggunanya. 
Kami juga ingin berterima kasih kepada Andreas Febrian, Lia Sadita, Fahrurrozi 
Rahman, Andre Tampubolon, dan Erik Dominikus atas kontribusinya dalam template 
ini. 

\vspace*{0.1cm}
\begin{flushright}
Depok, 27 Februari 2019\\[0.1cm]
\vspace*{1cm}
\penulis

\end{flushright}
%
%
\addChapter{LEMBAR PERSETUJUAN PUBLIKASI ILMIAH}
\input{src/00-front_matter/05-persetujuan_publikasi}
%
% 
\addChapter{ABSTRAK}
%
% Halaman Abstrak
%
% @author  Andreas Febrian
% @version 1.00
%

\chapter*{Abstrak}

\vspace*{0.2cm}
{
	\setlength{\parindent}{0pt}
	
	\begin{tabular}{@{}l l p{10cm}}
		Nama&: & \penulis \\
		Program Studi&: & \program \\
		Judul&: & \judul \\
		Pembimbing&: & \pembimbing \\
	\end{tabular}

	\bigskip
	\bigskip

	Tesis ini membahas kemampuan mahasiswa Fakultas Psikologi UI dalam mencari dan
	menggunakan informasi secara efektif dalam konteks \textit{active learning} dan \textit{self regulated
	learning} selama mereka mengikuti Program Pendidikan Dasar Pendidikan Tinggi.
	Penelitian ini adalah penelitian kualitatif dengan desain deskriptif. Hasil penelitian
	menyarankan bahwa perpustakaan perlu dilibatkan dalam pengembangan kurikulum;
	materi pendidikan pemakai perpustakaan harus dikembangkan sesuai dengan komponen-
	komponen yang ada dalam \textit{information literacy}; perpustakaan juga harus menyediakan
	sarana dan fasilitas yang mendukung peningkatan \textit{literacy} mahasiswa.

	\bigskip

	Kata kunci:\\
	Informasi, \textit{information literacy}, \textit{information skills}
}

\newpage
%
%
%
% Halaman Abstract
%
% @author  Andreas Febrian
% @version 1.00
%

\chapter*{Abstract}

\vspace*{0.2cm}
{
	\setlength{\parindent}{0pt}
	
	\begin{tabular}{@{}l l p{10cm}}
		Name&: & \penulis \\
		Study Program&: & \program \\
		Title&: & \judulInggris \\
		Counsellor&: & \pembimbing \\
	\end{tabular}

	\bigskip
	\bigskip

	The focus of this study is the freshman student of Faculty of Psychology at University of
	Indonesia experience of acquiring, evaluating and using information, when they enroll in
	“Program Dasar Pendidikan Tinggi (PDPT)”. The purpose of this study is to understand
	how freshman students acquire, evaluate and use information. Knowing this will allow
	library to identify changes should be made to improve user education program at
	University of Indonesia. This research is qualitative descriptive interpretive. The data
	were collected by means of deep interview. The researcher suggests that library should
	improve the user education program and provide facilities which can help students to be
	information literate.

	\bigskip

	Key words:\\
	Information literacy, information skills, information
}

\newpage

%
% Daftar isi, gambar, dan tabel
%
\phantomsection
\tableofcontents
\clearpage
\phantomsection
\listoffigures
\clearpage
\phantomsection
\listoftables
\clearpage
\phantomsection
\listoflistings
\clearpage

%
% Gunakan penomeran Arab (1, 2, 3, ...) setelah bagian ini.
%
\pagenumbering{arabic}

%
%
%
\input{src/01-body/01-bab1}
\input{src/01-body/02-bab2}
\input{src/01-body/03-bab3}
%-----------------------------------------------------------------------------%
\chapter{\babEmpat}
%-----------------------------------------------------------------------------%
\todo{tambahkan kata-kata pengantar bab 1 disini}

%-----------------------------------------------------------------------------%
\section{thesis.tex}
%-----------------------------------------------------------------------------%
Berkas ini berisi seluruh berkas Latex yang dibaca, jadi bisa dikatakan sebagai 
berkas utama. Dari berkas ini kita dapat mengatur bab apa saja yang ingin 
kita tampilkan dalam dokumen.


%-----------------------------------------------------------------------------%
\section{laporan\_setting.tex}
%-----------------------------------------------------------------------------%
Berkas ini berguna untuk mempermudah pembuatan beberapa template standar. 
Anda diminta untuk menuliskan judul laporan, nama, npm, dan hal-hal lain yang 
dibutuhkan untuk pembuatan template. 


%-----------------------------------------------------------------------------%
\section{istilah.tex}
%-----------------------------------------------------------------------------%
Berkas istilah digunakan untuk mencatat istilah-istilah yang digunakan. 
Fungsinya hanya untuk memudahkan penulisan.
Pada beberapa kasus, ada kata-kata yang harus selalu muncul dengan tercetak 
miring atau tercetak tebal. 
Dengan menjadikan kata-kata tersebut sebagai sebuah perintah \latex~tentu akan 
mempercepat dan mempermudah pengerjaan laporan. 


%-----------------------------------------------------------------------------%
\section{hype.indonesia.tex}
%-----------------------------------------------------------------------------%
Berkas ini berisi cara pemenggalan beberapa kata dalam bahasa Indonesia. 
\latex~memiliki algoritma untuk memenggal kata-kata sendiri, namun untuk 
beberapa kasus algoritma ini memenggal dengan cara yang salah. 
Untuk memperbaiki pemenggalan yang salah inilah cara pemenggalan yang benar 
ditulis dalam berkas hype.indonesia.tex.


%-----------------------------------------------------------------------------%
\section{pustaka.tex}
%-----------------------------------------------------------------------------%
Berkas pustaka.tex berisi seluruh daftar referensi yang digunakan dalam 
laporan. 
Anda bisa membuat model daftar referensi lain dengan menggunakan bibtex.
Untuk mempelajari bibtex lebih lanjut, silahkan buka 
\url{http://www.bibtex.org/Format}. 
Untuk merujuk pada salah satu referensi yang ada, gunakan perintah \bslash 
cite, e.g. \bslash cite\{lankton2008introduction\} yang akan akan memunculkan 
\cite{lankton2008introduction}


%-----------------------------------------------------------------------------%
\section{bab[1 - 6].tex}
%-----------------------------------------------------------------------------%
Berkas ini berisi isi laporan yang Anda tulis. 
Setiap nama berkas e.g. bab1.tex merepresentasikan bab dimana tulisan tersebut 
akan muncul. 
Sebagai contoh, kode dimana tulisan ini dibaut berada dalam berkas dengan nama 
bab4.tex. 
Ada enam buah berkas yang telah disiapkan untuk mengakomodir enam bab dari 
laporan Anda, diluar bab kesimpulan dan saran. 
Jika Anda tidak membutuhkan sebanyak itu, silahkan hapus kode dalam berkas 
thesis.tex yang memasukan berkas \latex~yang tidak dibutuhkan;  contohnya 
perintah \bslash include\{bab6.tex\} merupakan kode untuk memasukan berkas 
bab6.tex kedalam laporan.

%-----------------------------------------------------------------------------%
\section{Penulisan \textit{code} atau \textit{pseudocode} program}
%-----------------------------------------------------------------------------%

\subsection{\textit{Inline}}

Dengan perintah \verb|\verb|: \verb|System.out.println("Hello, World");| \\
Dengan perintah \textit{custom} \verb|\code|: \code{System.out.println("Hello, World"); }
Dengan perintah \verb|\mintinline|: \mintinline{java}{System.out.println("Hello, World"); }

\subsection{\textit{Multiline}}

Dengan perintah \verb|verbatim|: 

\begin{verbatim}	
public class HelloWorld {
    public static void main(String[] args) {
        // Prints "Hello, World" to the terminal window.
        System.out.println("Hello, World");
    }
}
\end{verbatim}

Dengan perintah \verb|minted|: Kode \ref{code:hw:minted}
\begin{listing}[H]
    \begin{minted}{python}
def binary_accuracy(y_true, y_pred):
    return K.mean(K.equal(y_true, K.round(y_pred)), axis=-1)


def categorical_accuracy(y_true, y_pred):
    return K.cast(K.equal(K.argmax(y_true, axis=-1),
                          K.argmax(y_pred, axis=-1)),
                  K.floatx())


def sparse_categorical_accuracy(y_true, y_pred):
    # reshape in case it's in shape (num_samples, 1) instead of (num_samples,)
    if K.ndim(y_true) == K.ndim(y_pred):
        y_true = K.squeeze(y_true, -1)
    # convert dense predictions to labels
    y_pred_labels = K.argmax(y_pred, axis=-1)
    y_pred_labels = K.cast(y_pred_labels, K.floatx())
    return K.cast(K.equal(y_true, y_pred_labels), K.floatx())


def top_k_categorical_accuracy(y_true, y_pred, k=5):
    return K.mean(K.in_top_k(y_pred, K.argmax(y_true, axis=-1), k), axis=-1)


def sparse_top_k_categorical_accuracy(y_true, y_pred, k=5):
    # If the shape of y_true is (num_samples, 1), flatten to (num_samples,)
    return K.mean(K.in_top_k(y_pred, K.cast(K.flatten(y_true), 'int32'), k),
                  axis=-1)
    \end{minted}
    \caption{An excerpt from keras: \url{https://github.com/keras-team/keras/blob/master/keras/metrics.py}}
    \label{code:hw:minted}
\end{listing}

Konfigurasi tampilan bisa dilakukan di \verb|uithesis.sty| dengan referensi dokumentasi di \url{https://github.com/gpoore/minted/blob/master/source/minted.pdf}
%-----------------------------------------------------------------------------%
\chapter{\babLima}
%-----------------------------------------------------------------------------%
\todo{Tambahkan kata-kata pengantar bab 5 disini.}


%-----------------------------------------------------------------------------%
\section{Mengubah Tampilan Teks}
%-----------------------------------------------------------------------------%
Beberapa perintah yang dapat digunakan untuk mengubah tampilan adalah: 
\begin{itemize}
	\item \bslash f \\
		Merupakan alias untuk perintah \bslash textit, contoh 
		\f{contoh hasil tulisan}.
	\item \bslash bi \\
		\bi{Contoh hasil tulisan}.
	\item \bslash bo \\
		\bo{Contoh hasil tulisan}.
	\item \bslash m \\
		Contoh\ hasil\ tulisan: $\alpha \not= \m{\alpha}$
	\item \bslash code \\ 
		\code{Contoh hasil tulisan}.
\end{itemize}


%-----------------------------------------------------------------------------%
\section{Memberikan Catatan}
%-----------------------------------------------------------------------------%
Ada dua perintah untuk memberikan catatan penulisan dalam dokumen yang Anda 
kerjakan, yaitu: 
\begin{itemize}
	\item \bslash todo \\
		Contoh: \\ \todo{Contoh bentuk todo.}
	\item \bslash todoCite \\ 
		Contoh: \todoCite
\end{itemize}


%-----------------------------------------------------------------------------%
\section{Menambah Isi Daftar Isi}
%-----------------------------------------------------------------------------%
Terkadang ada kebutuhan untuk memasukan kata-kata tertentu kedalam Daftar Isi. 
Perintah \bslash addChapter dapat digunakan untuk judul bab dalam Daftar isi. 
Contohnya dapat dilihat pada berkas thesis.tex.


%-----------------------------------------------------------------------------%
\section{Memasukan PDF}
%-----------------------------------------------------------------------------%
Untuk memasukan PDF dapat menggunakan perintah \bslash inpdf yang menerima satu 
buah argumen. Argumen ini berisi nama berkas yang akan digabungkan dalam 
laporan. PDF yang dimasukan degnan cara ini akan memiliki header dan footer 
seperti pada halaman lainnya. 

\inpdf{assets/pdfs/include}

Cara lain untuk memasukan PDF adalah dengan menggunakan perintah \bslash putpdf 
dengan satu argumen yang berisi nama berkas pdf. Berbeda dengan perintah 
sebelumnya, PDF yang dimasukan dengan cara ini tidak akan memiliki footer atau 
header seperti pada halaman lainnya. 

\putpdf{assets/pdfs/include}


%-----------------------------------------------------------------------------%
\section{Membuat Perintah Baru}
%-----------------------------------------------------------------------------%
Ada dua perintah yang dapat digunakan untuk membuat perintah baru, yaitu: 
\begin{itemize}
	\item \bslash Var \\
		Digunakan untuk membuat perintah baru, namun setiap kata yang diberikan
		akan diproses dahulu menjadi huruf kapital. 
		Contoh jika perintahnya adalah \bslash Var\{adalah\} makan ketika 
		perintah \bslash Var dipanggil, yang akan muncul adalah ADALAH. 
	\item \bslash var \\
		Digunakan untuk membuat perintah atau baru. 
\end{itemize}


\input{src/01-body/06-bab6}
\input{src/01-body/99-kesimpulan}

%
% Daftar Pustaka
%%
% Daftar Pustaka 
% 

% 
% Tambahkan pustaka yang digunakan setelah perintah berikut. 
% 
\begin{thebibliography}{4}

\bibitem{latex.intro}
{Clark, J. (n.d). \f{Introduction to LaTeX}.
26 Januari 2010. \url{http://frodo.elon.edu/tutorial/tutorial/node3.html}.}

\bibitem{chin2000learning}
{Chin, C., Brown, D. E. \f{Learning in science: A comparison of deep and surface approaches}.
2000. Journal of research in science teaching, 37(2), 109-138.}

\end{thebibliography}



% Alternatif manajemen daftar pustaka dengan \bibliography
\bibliographystyle{apalike}
\bibliography{pustaka}

%
% Lampiran 
%
% \begin{appendix}
% 	%
% @author  Andreas Febrian
% @version 1.00 
% 
% Hanya sebuah pembatas bertuliskan LAMPIRAN ditengah halaman. 
% 

\begin{titlepage}
	\centering 
	\vspace*{6cm}
	\noindent \Huge{LAMPIRAN}
	\addChapter{LAMPIRAN}
\end{titlepage}
% 	\setcounter{page}{2}\textbf{}
% 	%-----------------------------------------------------------------------------%
\addChapter{Lampiran 1}
\chapter*{Lampiran 1}
%-----------------------------------------------------------------------------%
% \end{appendix}

\end{document}